\documentclass[10pt]{article}
\usepackage[margin=0.9in]{geometry}
\usepackage{amsmath,amssymb,amsthm}
\usepackage[hidelinks]{hyperref}
\newtheorem{theorem}{Theorem}
\title{A counterexample to the three-shifted\\q-Motzkin divisibility conjecture}
\author{Principia Math In Progress}
\date{18 September 2026}
\begin{document}
\maketitle
\begin{abstract}
The size-five three-shifted Hankel determinant has value \(2\) at
\(q=-1\), while its conjectured divisor \(1+q^3\) vanishes.
This disproves the divisibility part of the intended conjecture.
The journal final remains unavailable for a corrections check;
this proof therefore carries zero new campaign credit.
\end{abstract}

\section*{Task 1: the original conjecture and a counterexample}
Write \(M_n(q)=\widetilde{\mathcal M}^{\dagger}_n(q)\) for the polynomials
defined on page 62 of \cite{OS} by
\[
 M_0(q)=1,\qquad
 M_{n+1}(q)=M_n(q)+\sum_{r=0}^{n-1}q^{r+1}M_r(q)M_{n-r-1}(q).
\]
Set
\[
 D_n(q)=\det[M_{i+j+3}(q)]_{i,j=0}^{n-1},
 \qquad t_n(q)=\sum_{r=0}^{\lfloor n/3\rfloor}q^{3r}.
\]
Conjecture A.3, reproduced as MathDB 339077, intends
\(D_n=(-1)^{\lfloor n/3\rfloor}q^{c_n}g_n(q)\), with
\(c_n\geq0\), \(g_n\in\mathbb Z_{\geq0}[q]\), and in particular
\(t_n\mid g_n\) when \(n\equiv2\pmod3\).
The printed display uses \(f_n\) where the ensuing text uses \(g_n\).
We resolve that notation in the intended way and disprove the
substantive divisibility assertion.

\begin{theorem}
We have \(D_5(-1)=2\). Consequently \(1+q^3\) does not divide \(D_5(q)\),
and no choice of monomial prefactor makes its normalized polynomial
divisible by \(t_5(q)=1+q^3\).
\end{theorem}
\begin{proof}
Put \(a_j=M_j(-1)\). Direct substitution in the defining recurrence gives
\[
 a_{n+1}=a_n+\sum_{r=0}^{n-1}(-1)^{r+1}a_ra_{n-r-1}.
\]
Starting at \(a_0=1\), its values through \(a_{11}\) are
\[
 (1,1,0,0,1,1,-1,-1,3,3,-6,-6).
\]
Thus \(D_5(-1)=\det A\), where
\[
 A=\begin{pmatrix}
 0&1&1&-1&-1\\
 1&1&-1&-1&3\\
 1&-1&-1&3&3\\
 -1&-1&3&3&-6\\
 -1&3&3&-6&-6
 \end{pmatrix}.
\]
Swap the first two rows, then eliminate the first two columns using
their unit pivots. The remaining block is
\[
 B=\begin{pmatrix}2&2&-2\\2&2&-3\\-2&-3&1\end{pmatrix},
 \qquad \det A=-\det B=2.
\]
Indeed, subtracting the first row of \(B\) from its second and adding
it to its third leaves lower rows \((0,0,-1)\), \((0,-1,-1)\);
hence \(\det B=-2\).
If \(t_5\mid g_5\) and \(D_5=-q^{c_5}g_5\), then evaluation at
\(-1\) would give \(D_5(-1)=0\), a contradiction.
\end{proof}

\subsection*{Scope and exact checks}
Evaluation at a negative value of \(q\) is legitimate: the claim is
polynomial divisibility, so a proposed identity must hold under every
specialization. The argument does not assume that \(q\) is a positive
probability parameter and does not exploit the \(f_n/g_n\) notation.
A single admissible size \(5\equiv2\pmod3\) disproves the universal claim.

As a supplemental certificate, exact polynomial elimination gives
\[
 D_5(q)=-q^{62}
 \bigl(1+2q+2q^3+2q^4+q^6+2q^7\bigr).
\]
Its remainder on division by \(1+q^3\) is the constant \(2\).
The elementary proof above needs only the twelve displayed integers
and the row operations, not this larger polynomial computation.

The verifier checks the integer recurrence, the determinant by all
\(120\) signed permutation products and by the displayed elimination,
the full polynomial determinant and its remainder, and twelve moment
polynomials by a separate weighted-walk computation. These are exact
checks of a finite counterexample, not numerical evidence for an
unbounded assertion.

All 68 pages of the current primary preprint, including the references,
were read. Its page 62 was also inspected visually. The 2019 journal
final \cite{OSfinal} was identified but could not be obtained in full;
possible corrections and the complete subsequent literature remain
an eligibility limitation. The earlier repository package for MathDB
339076 evaluates shift two and proposes shift three as a follow-up;
it does not supply this counterexample.
No independent peer review or proof-assistant certification is claimed.

\section*{Task 2: determine the correct divisors}
Classify the sizes \(n\equiv2\pmod3\) for which \(t_n\mid D_n\).
Compute the exact greatest common divisor with \(t_n\) in the remaining
cases, using roots of unity to isolate the obstructing factors.

\section*{Task 3: evaluate all three-shift determinants}
Find a uniform formula for \(D_n\), including its order at \(q=0\).
Separate the surviving coefficient-sign and square-factor assertions
from the divisibility assertion disproved here.

\section*{Task 4: root-of-unity specializations}
Find recurrences for \(M_n(\zeta)\) and \(D_n(\zeta)\) when \(\zeta\)
has finite multiplicative order. Determine whether the size-five
obstruction extends to an infinite family of sizes.

\section*{Task 5: higher shifts}
Investigate coefficient signs and cyclotomic factors of
\(\det[M_{i+j+s}(q)]\) as the shift \(s\) varies.
In particular, the separate all-even-shifts positivity conjecture,
MathDB 339078, is not settled by this counterexample.
Tasks 2--5 are proposed directions and carry no additional credit.

\begin{thebibliography}{9}
\bibitem{OS}
S.-j. Oh and T. Scrimshaw, \emph{Identities from representation theory},
arXiv:1805.00113v1 (2018), Appendix A, pp.~61--63.
\url{https://arxiv.org/abs/1805.00113v1}.
\bibitem{OSfinal}
S.-j. Oh and T. Scrimshaw, \emph{Identities from representation theory},
Discrete Mathematics \textbf{342}(9) (2019), 2493--2541.
\url{https://doi.org/10.1016/j.disc.2019.05.020}.
Bibliographic identification only; final full text not read.
\end{thebibliography}
\end{document}
